\section{Everything I did that we had not agreed, and what is left}\label{sec:extras}

You gave me leave to make choices during the night and asked me to list them.

\textbf{Changes to the plan.}
\begin{enumerate}
\item The learning rate, the steadied copy of the model, and the extra conditions on the
      ``flat'' ending (section \ref{sec:changes}). Two smaller details were settled when the
      programs were reviewed: on the finer measure a change of less than 0.002 bits per word
      counts as no change, and a run is not asked whether it is flat at the extra grading made
      when its time runs out.
\item The 32-bit round was run although the first round did not meet the rule (you had told
      me during the night to run everything). Under the
      agreed procedure it was run on 16 and 256 animals only, and on 256 animals only once:
      when a long run showed that its ``flat'' endings were wrong, I dropped the repeats still
      to come and used the time for long runs with the ``flat'' ending switched off, on 256
      and then on 1,024 animals (section \ref{sec:round32}).
\item The time limit was 30 minutes per run in the first round, 15 in the 32-bit round (40 for
      one run) and 5 in the round with unequal use, to fit everything into one night. Width 4,
      on a crowded machine, had 28 and 24 minutes.
\item Until half past one, models were trained one at a time. After that the long runs,
      width 4 and the small checks ran side by side on separate processor cores, which got
      more done before morning but made each run slower. So I give sets, not minutes,
      wherever it matters.
\item Width 4 was run as you asked, but alongside the long runs and not after them, each run
      on a single slow core. That is why its runs are short (section \ref{sec:four}).
\item Before the first run I set aside a copy of the programs that train and grade the
      models, and every run of the night used that copy. Later corrections to the working
      version change nothing in how models are trained or graded, and none reached the runs.
\end{enumerate}

\textbf{Extra analyses and checks.} Three needed no new training: the look inside the word
table (section \ref{sec:probe}), the check on whether models use the season word (section
\ref{sec:round32}), and the comparison after equal amounts of text (section \ref{sec:equal}). Two small checks did need training. One is the check on
the steadied copy (section \ref{sec:changes}). The other is the side check we had planned:
does it matter that every model carries rows for all 1,024 animals even when its language has
only 16? I built width 32 with rows for its 16 animals and no others. It learned them in 23, 23
and 23 sets; with the full table it had taken 23, 16 and 23. So the idle rows changed nothing.

\textbf{Not done.} Of the 32-bit round as planned (four languages, widths 8 to 64, three
repeats), the 64-animal language, width 64, width 8 on 1,024
animals and every third repeat were not run at all; table \ref{tab:large} shows what was run
on the large languages. The word-table check for width 4. A check of width 6, which I started
and gave up for lack of time. And grading at finer steps, which we had talked about.

\textbf{What I would do next,} in this order.
\begin{enumerate}
\item \emph{Settle width 8 on 256 animals with 32 bits each:} three repeats, trained ten times
      as long, with width 16 beside it. If it reaches 100 it was only slow. If it does not, ``too
      small'' becomes more likely, though a stand-still can never quite be told from a limit.
\item \emph{Test your prediction in a model where a new word does not bring its own storage.}
      Give each animal a row of random numbers that training may not change. The row is then
      only a name tag, and everything the model learns about an animal must be kept in the
      body, which has a fixed size and is shared by all words. Then add animals until the body
      is full. There your prediction could succeed or fail far more clearly than here.
\item \emph{Compare after equal amounts of text, as you suggested,} and drop the stopping test.
      And \emph{find the smallest width that works} between 4 and 8.
\end{enumerate}

\textbf{Where everything is.} The programs are in the folder \s{air} (\s{README.md} explains
each one), the runs in \s{air/runs}, \s{air/runs\_long} and \s{air/runs\_extra}, and the
trial-run logs in \s{air/logs}.
